\documentclass[12pt, a4paper]{article}

% Essential packages
\usepackage[utf8]{inputenc}
\usepackage{geometry}
\usepackage{setspace}
\usepackage{hyperref}
\usepackage{enumitem} % For customizing lists

% Page margins (Academic standard: 3cm left for binding, 2.5cm others)
\geometry{left=3cm, right=2.5cm, top=2.5cm, bottom=2.5cm}
\onehalfspacing     % Standard academic spacing

% Title Information
\title{\textbf{ECSP3000 Project Brief}\\[0.8em] \Large \textit{Designing and Implementing an Automated \& Reproducible Testbed for Linux Rootkit Detectors, Aiming to Identify and Implement a Novel Rootkit Technique}}
\author{Elliott Jelbert \\ 35949244 - ej1g24@soton.ac.uk \\ Supervisor: Dr. BooJoong Kang}

\begin{document}

\maketitle

\section*{1. Project Aims}

Current research into the evaluation of rootkits and rootkit detectors rely on bespoke evaluation methods %
which cannot be reliably reproduced by third partys. \\ \\ %
This tool aims to provide developers, researchers and system administrators with an automated and reproducible testbed for rootkits \& rootkit detectors on the modern Linux kernel to aid future development and testing.
\\ \\ Furthermore the testbed aims to provide greater transparency on what rootkit techniques remain evasive to standard detection tools.
\\ \\ As an extension, I will use this tool to identify and evaluate current avenues in the modern (7.2) %
Linux Kernel for rootkit design, finishing with an evaluation using the aforementioned tool of a new rootkit PoC.


\section*{2. Project Objectives}
This project will have the following aims which will be used to gauge success:
\begin{enumerate}
    \item To have a testing framework which, upon specifying a set of rootkits \& rootkit detectors, provides timely \& accurate results which can then be reproduced exactly upon rerun.
z    \item The system will aim to treat rootkits and rootkit detectors as pluggable modules, allowing developers to easily test new rootkit detectors or test against new rootkits.
    \item (Stretch) To research and develop a rootkit PoC for the modern 7.2 Linux Kernel, following with an evaluation using this testing framework.
    \item Using a dataset of results produced by the testbed or, if successful, the aforementioned PoC, to evaluate the current status of rootkit detection on the modern Linux Kernel.
\end{enumerate}

\section*{3. Limitations \& Scope}

In order to scope my project appropriately, here are a list of things that I currently plan on \textbf{not} doing.

\begin{enumerate}
    \item My testbed and therefore my rootkit research will only be concerned with Kernel Space Rootkits. This excludes user-space, firmware or any other rootkit classes from my research. %I am certain on this one
    \item Testing rootkits that implement anti-sandbox measures fall out of scope for this project due to the virtualised environment employed by the testing framework.
\end{enumerate}

% \section*{4. Technical \& Development Approach}

% Brief overview of how I will approach this project both technically and research wise

\end{document}